\ncdlesson{5}{Majas nye arbejde på caféen}{玛雅在咖啡馆的新工作}
{DU3 Modul 2（A2）｜对标 NCE 2 Lesson 11 \& Modul 2 \textit{Mundtlig Opgave 2 (Informationskløft)}}
{\textbf{间接疑问从句（Indirekte spørgsmål）}：是非问句转 \textit{om + 主语 + (ikke) + 动词} · 特殊问句转 \textit{hv-词 + 主语 + (ikke) + 动词}（\textbf{严禁倒装！}）· 第三人称信息差提问与转述}
{DU3 Modul 2 口试第二题（\textit{Informationskløft}）就是考你看着“玛雅在咖啡馆工作”这类图卡，用第三人称 \textit{hun/han} 向搭档连问 6 个聚焦问题！问句既可以用直接疑问句，也可以升级成礼貌的间接疑问句——但记住：\textbf{进了间接问句，动词千万别再倒装}！}

% ---------------------------------------------------------------------
\begin{foerduleaeser}
带着下面两个问题阅读 Tekst A 和 Tekst B：
\begin{enumerate}
  \item \textit{Hvor ligger den café, hvor Maja arbejder, og hvad skal hun lave om morgenen?}（玛雅工作的咖啡馆在哪里？她早上需要做哪些工作？）
  \item \textit{Hvor meget får Maja i løn i timen, og hvordan kommer hun på arbejde ifølge Tekst B?}（根据 Tekst B 的信息卡，玛雅的时薪是多少？她怎么去上班？）
\end{enumerate}
\end{foerduleaeser}

% ---------------------------------------------------------------------
\ncdsubsec{二、Tekst A \& Tekst B（双轨课文、发音提示与中文对照）}

\subsection*{Tekst A：\textit{Majas nye arbejde på caféen}（叙事短文 · 约 145 词）}

\begin{tipbox}{语篇说明}
结合 NCE 2 职场生活故事与 DU3 Modul 2 官方样题角色 Maja 编写。阅读时请特别注意\textbf{\color{NCDNavy}深蓝色标记的间接疑问词（\textit{om / hv-}）}后面紧跟着的\textbf{\color{NCDRed}红色“主语 + 动词”不倒装语序}！
\end{tipbox}

Maja har lige fået et nyt deltidsjob på en travl café ved havnen i Aarhus. I mandags var det hendes første arbejdsdag, og hendes venner fra sprogskolen var meget nysgerrige.

Før hun begyndte, spurgte hendes ven Kasper: ”\da{Ved du, hvornår} \dabold{du skal møde} om morgenen, og \da{om} \dabold{du også arbejder} i weekenden?” Maja smilte og svarede, at hun endnu ikke vidste, \da{hvor travlt} \dabold{der bliver} om lørdagen.

Nu kender hun hele planen. Hun møder klokken syv om morgenen for at bage friske boller, lave espresso og gøre bordene klar til kunderne. Hendes venner vil også gerne vide, \da{hvad} \dabold{hendes timeløn er}, og \da{om} \dabold{kollegerne er} søde. Maja fortæller glad, at hun tjener hundrede og femogtredive kroner i timen, og at alle på caféen hjælper hinanden.

Om morgenen ved hun stadig ikke, \da{hvilke kunder} \dabold{der kommer} ind først, men hun elsker allerede sit nye arbejde ved vandet.

\begin{transbox}
\small
\textbf{\color{NCDTeal}[Tekst A 中文精译]}\\
玛雅刚刚在奥胡斯港口边一家忙碌的咖啡馆找到了一份新的兼职工作。上周一是她第一天上班，语言学校的朋友们都非常好奇。

在开工前，她的朋友卡斯珀问她：“你知道你早上几点上班，还有你周末是不是也工作吗？”玛雅笑着回答说，她还不知道周六的时候会忙到什么程度。

现在她对整个排班已经了如指掌了。她早上七点到店，负责烘烤新鲜面包卷、冲意式浓缩咖啡以及把桌子收拾妥当迎接顾客。她的朋友们还很想知道她的时薪是多少，以及同事们好不好相处。玛雅开心地说，她每小时赚一百三十五克朗，而且咖啡馆里大家都会互相帮忙。

每天早晨，她依然不知道哪些客人会最先走进店里，但她已经深深爱上了这份海边的新工作。
\end{transbox}

\subsection*{Tekst B：\textit{DU3 Modul 2 Opgave 2 (Informationskløft)：Maja arbejder på en café}}

\textbf{考场形式说明}：两人一组抽卡完成\textbf{信息差任务（Informationskløft）}。考生 A 看到带有 6 项图片与关键词的提示卡（不能偷看 B 卡），必须用第三人称（\textit{hun/Maja}）发出 6 个精准提问；考生 B 根据自己手里 B 卡的图片信息作答（\textbf{仅提问方在此题计分，4--6 分及格}）。

\begin{tcolorbox}[enhanced, breakable, colback=NCDBgGray, colframe=NCDGray, coltitle=white, fonttitle=\bfseries, title={Modul 2 口试信息差 6 张提示卡全示范（Kursist A 提问双轨模板 + Kursist B 完整回答）}, arc=2pt, boxrule=0.8pt, borderline west={3pt}{0pt}{NCDTeal}, left=8pt, right=8pt, top=7pt, bottom=7pt]
\renewcommand{\arraystretch}{1.28}
{\small
\begin{tabularx}{\linewidth}{@{} >{\bfseries\raggedright\arraybackslash}p{1.9cm} >{\raggedright\arraybackslash}p{4.0cm} >{\raggedright\arraybackslash}p{4.7cm} >{\raggedright\arraybackslash}X @{}}
\toprule
\textbf{提示词} & \textbf{直接问法（基础满分）} & \textbf{礼貌间接问法（高分秀语序）} & \textbf{Kursist B 回答（第三人称）} \\
\midrule
1. Arbejdssted\newline （地点） & \textit{Hvor arbejder Maja henne?} & \textit{Kan du fortælle mig, \da{hvor} \dabold{Maja arbejder}?} & \textit{Hun arbejder på Café Havnen i Aarhus.} \\
\midrule
2. Arbejdstid\newline （工时） & \textit{Hvornår møder hun, og arbejder hun i weekenden?} & \textit{Jeg vil gerne vide, \da{hvornår} \dabold{hun møder}, og \da{om} \dabold{hun arbejder} i weekenden.} & \textit{Hun møder kl. 7 om morgenen, og hun arbejder hver anden lørdag.} \\
\midrule
3. Løn\newline （薪水） & \textit{Hvor meget får hun i løn i timen?} & \textit{Ved du, \da{hvor meget} \dabold{hun tjener} i timen?} & \textit{Hun tjener 135 kroner i timen.} \\
\midrule
4. Opgaver\newline （任务） & \textit{Hvad er hendes arbejdsopgaver?} & \textit{Kan du sige mig, \da{hvad} \dabold{hun skal lave} på caféen?} & \textit{Hun bager boller, laver kaffe og betjener kunderne.} \\
\midrule
5. Kolleger\newline （同事） & \textit{Har hun mange kolleger, og er de søde?} & \textit{Jeg vil gerne vide, \da{om} \dabold{hun har} søde kolleger.} & \textit{Ja, hun har fire kolleger, og de er rigtig søde og hjælpsomme.} \\
\midrule
6. Transport\newline （交通） & \textit{Hvordan kommer hun på arbejde?} & \textit{Kan du fortælle mig, \da{hvordan} \dabold{hun kommer} hen til caféen?} & \textit{Hun cykler langs havnen, og det tager 15 minutter.} \\
\bottomrule
\end{tabularx}
}
\end{tcolorbox}

\begin{tipbox}{本课发音与吞音要点（Udtale \& Reduktioner）}
\begin{itemize}
  \item \textbf{咖啡馆与工作高频词}：\textit{caféen} \ipa{kʰaˈfeːˀən}（重音在 \textit{-fe-}）；\textit{arbejdstid} \ipa{ˈɑːbɑi̯ˀdsˌtˢiðˀ}（前面 \textit{d} 偏硬，后面 \textit{-tid} 为软 d + stød）；\textit{deltidsjob} \ipa{ˈdeːlˌtˢiðsˌjʌb}。
  \item \textbf{薪酬与同事词的静音字母}：\textit{løn} \ipa{ˈlønˀ}（短圆唇音 + stød）；\textit{kolleger} \ipa{kʰoˈleːɐ}（重音在 \textit{-le-}，\textbf{中间的 g 不发音}！）；\textit{opgaver} \ipa{ˈʌbˌɡæːvɐ}。
  \item \textbf{间接疑问词与吞音连读}：\textit{hvornår} \ipa{vɒˈnɒːˀ}（\textbf{词首 h 不发音}）；\textit{hvordan} \ipa{vɒˈdan}；\textit{hvor meget} \ipa{vɒ ˈmɑi̯əð}（\textit{g} 滑化为 \ipa{i̯}，\textit{-et} 读 \ipa{əð}）；\textit{om hun} \ipa{ʌm hun}（轻读带过）。
\end{itemize}
\end{tipbox}

\begin{transbox}
\small
\textbf{\color{NCDTeal}[Tekst B 中文对照]}\\
1. \textbf{工作地点}：玛雅在哪里工作？/ 你能告诉我玛雅在哪里工作吗？—— 她在奥胡斯的港口咖啡馆工作。\\
2. \textbf{工作时间}：她几点上班，周末工作吗？/ 我很想知道她几点上班，以及她周末是否工作。—— 她早上 7 点上班，隔周六上一次班。\\
3. \textbf{薪水}：她每小时挣多少钱？/ 你知道她时薪多少吗？—— 她每小时挣 135 克朗。\\
4. \textbf{工作任务}：她的工作任务是什么？/ 你能跟我说说她在咖啡馆要做什么吗？—— 她烤面包卷、做咖啡并服务顾客。\\
5. \textbf{同事}：她同事多吗，人好不好？/ 我想知道她有没有友善的同事。—— 是的，她有四位同事，大家都非常友善又乐于助人。\\
6. \textbf{交通方式}：她怎么去上班？/ 你能告诉我她怎么去咖啡馆吗？—— 她沿着港口骑自行车，大概要 15 分钟。
\end{transbox}

% ---------------------------------------------------------------------
\ncdsubsec{三、Gloser（核心生词表）}

\begin{tipbox}{词表重点}
\textit{en kollega} 的复数有两种官方拼写：\textit{kolleger} 或 \textit{kollegaer}（两者完全正确）；\textit{en løn}（工资）通常只有单数，不可数；\textit{at vide} 是极高频不规则动词（\textit{ved, vidste, har vidst}，注意过去式 \textit{vidste} 里的 d 不发音！）。
\end{tipbox}

\begin{glosertable}
连词 & om & \textit{(引导间接是非疑问从句)} & 是否（绝不可用 \textit{hvis} 代替！） \\
\midrule
动词 & at vide & \textit{(ved, vidste, har vidst)} & 知道，了解 \\
动词 & at spørge & \textit{(spørger, spurgte, har spurgt)} & 询问，提问（\textit{spørge om} 问关于……） \\
动词 & at fortælle & \textit{(fortæller, fortalte, har fortalt)} & 讲述，告诉 \\
动词 & at møde & \textit{(møder, mødte, har mødt)} & （按时）到岗上班；遇见 \\
动词 & at tjene & \textit{(tjener, tjente, har tjent)} & 挣（钱），赚取 \\
动词 & at bage & \textit{(bager, bagte, har bagt)} & 烘烤，烘焙 \\
动词 & at betjene & \textit{(betjener, betjente, har betjent)} & 服务，接待（顾客） \\
动词 \dabold{[er]} & at cykle & \textit{(cykler, cyklede, \dabold{er}/har cyklet)} & 骑自行车（有终点用 \textit{er}） \\
\midrule
名词 & en arbejdsplads / et arbejdssted & \textit{(-pladsen, -er / -stedet, -steder)} & 工作场所，工作地点 \\
名词 & en arbejdstid & \textit{(arbejdstiden, arbejdstider, -ne)} & 工作时间，班次 \\
名词 & en timeløn & \textit{(timelønnen, kun ental)} & 时薪（\textit{i timen} 每小时） \\
名词 & en arbejdsopgave & \textit{(arbejdsopgaven, -r, -rne)} & 工作任务 \\
名词 & en kollega & \textit{(kollegaen, kolleger/kollegaer, -ne)} & 同事 \\
名词 & en kunde & \textit{(kunden, kunder, kunderne)} & 顾客，客户 \\
名词 & en bolle & \textit{(bollen, boller, bollerne)} & 小圆面包卷 \\
形容词 & nysgerrig / hjælpsom & \textit{nysgerrig (-t, -e)} / \textit{hjælpsom (-t, -me)} & 好奇的 / 乐于助人的 \\
\end{glosertable}

% ---------------------------------------------------------------------
\ncdsubsec{四、Grammatik \& Skriveteknik（核心语法与 Modul 2 应试技巧精讲）}

\begin{grammatikbox}{语法点 1：间接疑问从句（Indirekte spørgsmål）三大铁律}
把一个直接问句放进 \textit{Jeg ved ikke... / Kan du fortælle mig... / De spørger...} 后面，它就变成了一个\textbf{从句（Ledsætning）}！90\% 的初学者都会错在“忘了取消倒装”：

\vspace{0.3em}
\renewcommand{\arraystretch}{1.28}
{\small
\begin{tabularx}{\linewidth}{@{} >{\bfseries\raggedright\arraybackslash}p{2.3cm} >{\raggedright\arraybackslash}p{4.3cm} >{\raggedright\arraybackslash}p{4.2cm} >{\raggedright\arraybackslash}X @{}}
\toprule
\textbf{原问句类型} & \textbf{直接问句（动词在前）} & \textbf{间接疑问从句公式} & \textbf{正确例句（严禁倒装！）} \\
\midrule
\textbf{1. 是非问句}\newline (Ja/Nej) & \textit{Arbejder hun i weekenden?}\newline \textit{Kommer han ikke i dag?} & \da{om} + 主语 + \dabold{(ikke)} + \textbf{动词} & \textit{Jeg ved ikke, \da{om} \dabold{hun arbejder} i weekenden.}\newline \textit{Ved du, \da{om} han \dabold{ikke kommer}?} \\
\midrule
\textbf{2. 普通 \textit{hv-} 问句} & \textit{Hvornår møder Maja?}\newline \textit{Hvorfor kom hun ikke?} & \da{hv-词} + 主语 + \dabold{(ikke)} + \textbf{动词} & \textit{Kan du sige mig, \da{hvornår} \dabold{Maja møder}?}\newline \textit{Jeg ved ikke, \da{hvorfor} hun \dabold{ikke kom}.} \\
\midrule
\textbf{3. \textit{hvem/hvad} 作主语} & \textit{Hvem kommer først?}\newline \textit{Hvad sker der her?} & \da{hvem/hvad} + \dabold{der} + \textbf{动词} & \textit{Hun ved ikke, \da{hvem} \dabold{der} kommer først.}\newline \textit{Fortæl mig, \da{hvad} \dabold{der} sker!} \\
\bottomrule
\end{tabularx}
}
\vspace{0.3em}

\textbf{两大致命禁区}：
\begin{enumerate}
  \item \textbf{严禁保留倒装}：绝不能说 *\textit{Kan du fortælle mig, hvornår \dabold{møder Maja}?} $\rightarrow$ 必须恢复“主语 + 动词”：*\textit{...hvornår \dabold{Maja møder}?}
  \item \textbf{\textit{om} 绝不能换成 \textit{hvis}}：英语的 *if* 既表示“如果”又表示“是否”，但丹麦语严格分开：“如果”用 \textit{hvis}，“是否”必须用 \dabold{om}！
\end{enumerate}
\end{grammatikbox}

\begin{grammatikbox}{语法点 2：DU3 Modul 2 口语 Opgave 2（Informationskløft）考场拿满分法则}
根据丹麦官方评分细则，Opgave 2 共 6 张图片/提示词，每个合格问题得 1 分（4 分及格）：
\begin{itemize}
  \item \textbf{人称必须用第三人称（\textit{hun / han / [Navn]}）}：卡上是别人的故事（如 Maja / Kasper），\textbf{不要对着考官或搭档问 *du*}（除非你说 *Kan \textbf{du} fortælle mig, hvor \textbf{hun} arbejder?*）。
  \item \textbf{不要给右上角头像提问}：头像只用来告诉你这个人物的名字和性别，6 个问题必须对准 6 个分类框里的提示词和图片！
  \item \textbf{安全稳分策略}：6 个问题里，可以用 3--4 个清晰简洁的直接问句（如 \textit{Hvornår møder hun?}）+ 2--3 个礼貌间接问句（如 \textit{Jeg vil gerne vide, om hun tjener godt.}）来展示 A2+ 的语法掌控力！
\end{itemize}
\end{grammatikbox}

% ---------------------------------------------------------------------
\ncdsubsec{五、Mønsterøvelser（关键句型操练）}

\textbf{练习 A：将下列直接问句改写为间接疑问句（注意取消倒装、用 \textit{om} 以及 \textit{ikke/der} 的位置！）}
\begin{quote}
\small 例：\textit{Hvornår møder Maja?} $\rightarrow$ \textit{Jeg ved ikke, hvornår Maja møder.}
\end{quote}
\begin{enumerate}
  \item \textit{Arbejder Maja også om søndagen?} $\rightarrow$ \textit{Kan du fortælle mig, ...}
  \item \textit{Hvor meget tjener hun i timen?} $\rightarrow$ \textit{Jeg vil gerne vide, ...}
  \item \textit{Hvordan kommer hun på arbejde hver morgen?} $\rightarrow$ \textit{Ved du, ...}
  \item \textit{Hvorfor møder hun ikke klokken otte i dag?} $\rightarrow$ \textit{Jeg ved ikke, ...}
  \item \textit{Er hendes kolleger søde og hjælpsomme?} $\rightarrow$ \textit{De spørger, ...}
  \item \textit{Hvem bager bollerne på caféen?} $\rightarrow$ \textit{Kan du sige mig, ...}（提示：\textit{hvem} 作主语！）
  \item \textit{Hvilke opgaver har hun om eftermiddagen?} $\rightarrow$ \textit{Jeg vil gerne vide, ...}
\end{enumerate}

\textbf{练习 B：火眼金睛纠错（找出并改正下列间接疑问句中的语序或选词错误）}
\begin{enumerate}
  \item \textit{Jeg ved ikke, hvor arbejder Maja henne.} $\rightarrow$
  \item \textit{Kan du fortælle mig, hvis hun har fri i weekenden?} $\rightarrow$
  \item \textit{De vil gerne vide, hvorfor kan hun ikke komme i aften.} $\rightarrow$
  \item \textit{Jeg vil gerne høre, hvad laver hun på caféen.} $\rightarrow$
  \item \textit{Hun spørger, om han arbejder aldrig om aftenen.} $\rightarrow$
  \item \textit{Jeg ved ikke, hvem kommer ind i butikken.} $\rightarrow$
\end{enumerate}

\textbf{练习 C：中译丹实战（口语信息差提问与转述）}
\begin{enumerate}
  \item 你能告诉我玛雅在哪家咖啡馆工作吗？
  \item 我想知道她周末是不是也上班。
  \item 他知道玛雅一小时挣多少钱吗？
  \item 我不知道为什么公交车今天没有按时来。
  \item 她朋友问她有哪些同事，以及她怎么去上班。
\end{enumerate}

% ---------------------------------------------------------------------
\ncdsubsec{六、Modul 2 实战演练与参考答案}

\begin{tcolorbox}[enhanced, breakable, colback=NCDBgGray, colframe=NCDGray, coltitle=white, fonttitle=\bfseries, title={Mundtlig Opgave 2（DU3 Modul 2 口语信息差模拟卡）：Kaspers job i et supermarked}, arc=2pt, boxrule=0.8pt, borderline west={3pt}{0pt}{NCDRed}, left=10pt, right=10pt, top=7pt, bottom=7pt]
\textbf{任务一（提问练习）}：假设你的搭档手里有关于 \textbf{Kasper} 在超市工作的 B 卡，你只有下面 6 个提示词，请用\textbf{第三人称 \textit{han / Kasper}} 提出 6 个聚焦问题（至少 3 个使用间接疑问句！）：
\begin{center}
\small
\textbf{1. sted / by} \quad|\quad \textbf{2. arbejdstider} \quad|\quad \textbf{3. løn} \quad|\quad \textbf{4. arbejdsopgaver} \quad|\quad \textbf{5. kolleger / pause} \quad|\quad \textbf{6. transport til arbejde}
\end{center}
\textbf{任务二（口语转述）}：收集完信息后，请根据下面的参考信息（奥胡斯 Netto 超市，每周三、五下午及周六，130 kr./小时，摆货与收银，3 位好同事，坐轻轨 Letbanen 20 分钟），用 85--105 词写出一段连贯的第三人称介绍！
\end{tcolorbox}

\begin{answerbox}
\textbf{Før du læser 参考回答}
\begin{enumerate}
  \item \textit{Caféen ligger ved havnen i Aarhus. Om morgenen møder Maja kl. 7 for at bage friske boller, lave espresso og gøre bordene klar til kunderne.}
  \item \textit{Maja får 135 kroner i timen, og hun cykler langs havnen på arbejde (det tager 15 minutter).}
\end{enumerate}

\textbf{练习 A 答案（间接疑问句改写）}
\begin{enumerate}
  \item \textit{Kan du fortælle mig, \da{om} Maja også arbejder om søndagen?}
  \item \textit{Jeg vil gerne vide, \da{hvor meget} hun tjener i timen.}
  \item \textit{Ved du, \da{hvordan} hun kommer på arbejde hver morgen?}
  \item \textit{Jeg ved ikke, \da{hvorfor} hun \dabold{ikke møder} klokken otte i dag.}
  \item \textit{De spørger, \da{om} hendes kolleger er søde og hjælpsomme.}
  \item \textit{Kan du sige mig, \da{hvem der} bager bollerne på caféen?}（\textit{hvem} 作主语必须加 \textit{der}！）
  \item \textit{Jeg vil gerne vide, \da{hvilke opgaver} hun har om eftermiddagen.}
\end{enumerate}

\textbf{练习 B 答案（纠错）}
\begin{enumerate}
  \item 取消倒装：\textit{Jeg ved ikke, hvor \dabold{Maja arbejder} henne.}
  \item “是否”用 \textit{om} 不用 \textit{hvis}：\textit{Kan du fortælle mig, \dabold{om} hun har fri i weekenden?}
  \item 取消倒装且 \textit{ikke} 前置：\textit{De vil gerne vide, hvorfor \dabold{hun ikke kan} komme i aften.}
  \item 取消倒装：\textit{Jeg vil gerne høre, hvad \dabold{hun laver} på caféen.}
  \item 副词 \textit{aldrig} 在动词前：\textit{Hun spørger, om han \dabold{aldrig arbejder} om aftenen.}
  \item \textit{hvem} 作主语加 \textit{der}：\textit{Jeg ved ikke, hvem \dabold{der} kommer ind i butikken.}
\end{enumerate}

\textbf{练习 C 答案}
\begin{enumerate}
  \item \textit{Kan du fortælle mig, hvilken café Maja arbejder på?}
  \item \textit{Jeg vil gerne vide, om hun også arbejder i weekenden.}
  \item \textit{Ved han, hvor meget Maja tjener i timen?}
  \item \textit{Jeg ved ikke, hvorfor bussen ikke kom til tiden i dag.}
  \item \textit{Hendes venner spørger, hvilke kolleger hun har, og hvordan hun kommer på arbejde.}
\end{enumerate}

\textbf{实战任务一：Kasper 6 题提问参考（直接 + 间接混合示范）}
\begin{enumerate}
  \item \textit{Kan du fortælle mig, \da{hvor} Kasper arbejder henne?}
  \item \textit{Jeg vil gerne vide, hvornår han arbejder, og \da{om} han også arbejder i weekenden.}
  \item \textit{Ved du, \da{hvor meget} han får i løn i timen?}
  \item \textit{Hvad er hans arbejdsopgaver i supermarkedet?}
  \item \textit{Kan du sige mig, \da{hvordan} hans kolleger er, og \da{om} de holder pause sammen?}
  \item \textit{Hvordan kommer han hen til sit arbejde?}
\end{enumerate}

\textbf{实战任务二：第三人称口语转述范文（88 词）}\\
\textit{Kasper har fået et nyt deltidsjob i et supermarked, som ligger i Aarhus centrum. Jeg spurgte ham i går, \da{hvornår} han arbejder, og han fortalte, at han har vagter onsdag og fredag eftermiddag samt hver lørdag. Han tjener hundrede og tredive kroner i timen. Hans vigtigste arbejdsopgaver er at sætte varer på hylderne og sidde ved kassen. Heldigvis ved jeg nu, \da{at} han har tre rigtig søde kolleger, og at de altid spiser frokost sammen i pausen. Han tager letbanen til arbejde, og turen tager kun tyve minutter.}

\small\textbf{点评}：清晰展示 Opgave 2 的“第三人称聚焦提问法”与“信息收集后转述”。示范问句交替使用了直接问句与带 \textit{om / hv-} 的间接疑问句（严格做到主语在动词前），88 词转述文本流畅串联地点、工时、薪资、任务、同事与交通全部 6 个节点。
\end{answerbox}
